\thispagestyle{empty}
\begin{center}
	
	\large
	
	\textbf{JOSÉ FERNANDO MORALES DE LEÓN}
	
	\vfill
	
	\begin{figure}[H]
		\centering
		\includegraphics[width=0.3\textwidth]{logoutp.png}
	\end{figure}

	\vfill
	
	\textbf{DESARROLLO DE UN PROTOTIPO DE UN DISPOSITIVO DE
	MEDICIÓN DE TEMPERATURA PARA LA EMPRESA SIGCSA}
	
	\vfill
	
	\textbf{Universidad Tecnológica de Panamá \\
	Facultad de Ingeniería en Sistemas Computacionales \\
	Licenciatura en Ingeniería de Sistemas y Computación}
	
	\vfill
	
	\textbf{Panamá\\
	2018}
		
\end{center}
\clearpage

\thispagestyle{empty}

\begin{center}

\large
JOSÉ FERNANDO MORALES DE LEÓN

\vfill

DESARROLLO DE UN PROTOTIPO DE UN DISPOSITIVO DE
MEDICIÓN DE TEMPERATURA PARA LA EMPRESA SIGCSA

\vfill

Tesis presentada a la Facultad de Ingeniería en Sistemas
Computacionales de la Universidad Tecnológica de Panamá
para la obtencion del Título de

\vfill

Licenciatura en Ingeniería de Sistemas y Computación

\vfill

Asesora:\\
M.Sc. Doris Elena Gutiérrez Rosales

\vfill
Panamá\\
2018

\end{center}

\clearpage

\thispagestyle{plain}

\begin{flushright}
	\vspace*{\fill}
	
	\textit{Dedico esto a mi familia, especialmente a mi abuela Emma De León Cabrera, gracias por su apoyo incondicional.}
	
	\vspace*{\fill}
\end{flushright}

\clearpage

\thispagestyle{plain}
\begin{center}
	\huge
	\textbf{Agradecimientos}
\end{center}

\par \noindent
Agradezco a Dios y a mi madre.

\clearpage

\thispagestyle{plain}
\begin{center}
	\huge
	\textbf{Resumen}
\end{center}

\par \noindent
La presente tesis presenta el diseño y desarrollo de un prototipo para la captura de medidas de temperatura. El proyecto busca reemplazar la actividad de capturar de datos por una persona y dejarle el trabajo meramente al dispositivo y su aplicación. La propuesta fue desarrollada luego de realizar una investigación del proceso actual que ejecuta la empresa Soluciones Integrales en Gestión de Calidad, S.A (SIGCSA) para la caracterización de medios isotermos, los estándares implementados en dicho proceso y recomendaciones del personal de la empresa.\\ 
El resultado final de este trabajo permitirá al personal de SIGCSA realizar otros labores; mientras su teléfono capturará las mediciones de temperaturas de los equipos a caracterizar, disminuyendo el tiempo de servicio técnico manteniendo la calidad de servicio a sus clientes.

\par \noindent
\textit{\textbf{Palabras Claves:} prototipo, caracterización, medios isotermos; mediciones de temperatura}

\clearpage
